\begin{frame}{1.3 Transferência paralela}
  Vários bits são enviados ao mesmo tempo, cada um por uma linha.
  \begin{center}
  \begin{tikzpicture}[every node/.style={align=center}, >=Latex, node distance=2cm]
    \node[draw, minimum width=2cm, minimum height=2.4cm] (a) {Origem};
    \node[draw, minimum width=2cm, minimum height=2.4cm, right=4cm of a] (b) {Destino};
    \foreach \y/\n in {0.8/0,0.25/1,-0.3/2,-0.85/3} {\draw[->, thick] (a.east |- 0,\y) -- node[above, font=\scriptsize] {$D_\n$} (b.west |- 0,\y);}
  \end{tikzpicture}
  \end{center}
  \begin{itemize}
    \item Exemplo: barramento interno de 8, 16 ou 32 bits.
    \item Vantagem: muitos bits por ciclo.
    \item Desvantagens: mais fios, conectores maiores e risco de \textbf{skew} (bits chegando em tempos diferentes).
  \end{itemize}
\end{frame}
